\chapter{Square Roots}\label{ch:g9:sqrt}

How long is the diagonal of a square of side $1$? The Pythagorean theorem
answers $\sqrt 2$, a number whose square is $2$ --- and which turns out
not to be a \omterm{def:g9:fractions:fraction}{fraction}. This chapter defines \omterm{def:g9:sqrt:def}{square roots}, establishes the
rules for computing with them, and teaches how to simplify expressions
like $\sqrt{75}$.

\section{Definition and first properties}

\begin{definition}[Square root]\label{def:g9:sqrt:def}
Let $a \geq 0$. The \emph{square root}\index{square root} of $a$, written
$\sqrt a$, is the unique \emph{nonnegative} number whose square is $a$:
\[
\sqrt a \geq 0
\qquad\text{and}\qquad
\left(\sqrt a\right)^2 = a .
\]
Negative numbers have no square root, since every square is nonnegative.
\end{definition}

\begin{example}
$\sqrt{49} = 7$, $\sqrt 0 = 0$, $\sqrt 1 = 1$,
$\sqrt{2.25} = 1.5$. The first \omterm{def:g9:sqrt:def}{square roots} to know by heart are those of
the \emph{perfect squares}: $1, 4, 9, 16, 25, 36, 49, 64, 81, 100, 121,
144$.
\end{example}

\begin{omfigure}
\begin{tikzpicture}[scale=1.35]
  \draw[thick, fill=omDef!8] (0,0) rectangle (2,2);
  \draw[omThm, thick] (0,0) -- (2,2);
  \node[below, font=\small] at (1,0) {$1$};
  \node[left, font=\small] at (0,1) {$1$};
  \node[omThm, above left=-2pt, font=\small, rotate=45] at (1.08,0.92)
    {$\sqrt2$};
  \draw[thin] (1.75,0) -- (1.75,0.25) -- (2,0.25);
\end{tikzpicture}

{\small The diagonal of a unit square has length $\sqrt{1^2 + 1^2} =
\sqrt2 \approx 1.414$, by the Pythagorean theorem
(\cref{ch:g9:trig} recalls it). This number is irrational: its decimals
never repeat.}
\end{omfigure}

\begin{proposition}[Square root of a square]\label{prop:g9:sqrt:square}
For every real number $x$ (positive or not):
\[
\sqrt{x^2} = \abs{x} =
\begin{cases}
x & \text{if } x \geq 0,\\
-x & \text{if } x < 0 .
\end{cases}
\]
\end{proposition}

\begin{proof}
The number $\abs{x}$ is nonnegative and its square is $x^2$ (a number and
its opposite have the same square). Being the unique nonnegative number
of square $x^2$, it is $\sqrt{x^2}$.
\end{proof}

\begin{example}
$\sqrt{(-5)^2} = \sqrt{25} = 5 = \abs{-5}$: the \omterm{def:g9:sqrt:def}{square root} ``forgets''
the sign, it does not restore it. In particular $\sqrt{x^2} = x$ is
\emph{wrong} for negative $x$.
\end{example}

\section{Products and quotients of square roots}

\begin{theorem}[Multiplication and division rules]\label{thm:g9:sqrt:rules}
For all $a \geq 0$ and $b \geq 0$:
\[
\sqrt{a b} = \sqrt a \times \sqrt b,
\qquad\text{and for } b > 0:\quad
\sqrt{\frac ab} = \frac{\sqrt a}{\sqrt b} .
\]
\end{theorem}

\begin{proof}
The number $\sqrt a \times \sqrt b$ is nonnegative (\omterm{def:g2:mult:def}{product} of
nonnegative numbers), and its square is
\[
\left(\sqrt a \times \sqrt b\right)^2
= \left(\sqrt a\right)^2 \times \left(\sqrt b\right)^2
= a b .
\]
By uniqueness of the nonnegative number whose square is $ab$, it equals
$\sqrt{ab}$. The \omterm{def:g3:division:remainder}{quotient} rule is proved the same way.
\end{proof}

\begin{remark}[No such rule for sums!]
$\sqrt{a + b}$ is \emph{not} $\sqrt a + \sqrt b$ in general:
\[
\sqrt{9 + 16} = \sqrt{25} = 5,
\qquad\text{but}\qquad
\sqrt 9 + \sqrt{16} = 3 + 4 = 7 .
\]
\end{remark}

\begin{method}[Simplifying $\sqrt n$]\label{met:g9:sqrt:simplify}
To simplify the \omterm{def:g9:sqrt:def}{square root} of an integer:
\begin{enumerate}
  \item find the largest perfect square dividing $n$ (factor $n$ if
        needed);
  \item split with the \omterm{def:g2:mult:def}{product} rule and extract the square:
        $\sqrt{k^2 m} = k\sqrt m$;
  \item check that no perfect square \omterm{def:g9:arith:divisor}{divides} what is left under the
        root.
\end{enumerate}
\end{method}

\begin{example}\label{ex:g9:sqrt:simplify}
Simplify $\sqrt{75}$: since $75 = 25 \times 3$,
\[
\sqrt{75} = \sqrt{25 \times 3} = \sqrt{25} \times \sqrt 3 = 5\sqrt3 .
\]
Simplify $\sqrt{72}$: since $72 = 36 \times 2$,
$\sqrt{72} = 6\sqrt2$. And a \omterm{def:g3:division:remainder}{quotient}:
$\sqrt{\dfrac{49}{16}} = \dfrac{\sqrt{49}}{\sqrt{16}} = \dfrac74$.
\end{example}

\begin{example}[Adding square roots]\label{ex:g9:sqrt:add}
\omterm{def:g1:addition:def}{Sums} of \omterm{def:g9:sqrt:def}{square roots} simplify only when the roots are \emph{alike}.
Compute $\sqrt{12} + \sqrt{27} - \sqrt{48}$, simplifying each term first:
\begin{align*}
\sqrt{12} &= \sqrt{4 \times 3} = 2\sqrt3, \\
\sqrt{27} &= \sqrt{9 \times 3} = 3\sqrt3, \\
\sqrt{48} &= \sqrt{16 \times 3} = 4\sqrt3,
\end{align*}
so the \omterm{def:g1:addition:def}{sum} is $2\sqrt3 + 3\sqrt3 - 4\sqrt3 = (2 + 3 - 4)\sqrt3 = \sqrt3$.
\end{example}

\begin{example}[Expanding with square roots]\label{ex:g9:sqrt:expand}
The identities of algebra apply to \omterm{def:g9:sqrt:def}{square roots}. Expand
$\left(\sqrt5 + 2\right)^2$ with $(a+b)^2 = a^2 + 2ab + b^2$:
\[
\left(\sqrt5 + 2\right)^2 = 5 + 2 \times 2\sqrt5 + 4 = 9 + 4\sqrt5 .
\]
And with the third identity, $(a+b)(a-b) = a^2 - b^2$:
\[
\left(\sqrt7 + \sqrt3\right)\left(\sqrt7 - \sqrt3\right) = 7 - 3 = 4 :
\]
the \omterm{def:g2:mult:def}{product} of two irrational numbers can be an integer.
\end{example}

\section{The equation \texorpdfstring{$x^2 = a$}{x2 = a}}

\begin{theorem}[Solving $x^2 = a$]\label{thm:g9:sqrt:equation}
\begin{enumerate}
  \item If $a > 0$, the \omterm{def:g8:equations:def}{equation} $x^2 = a$ has exactly two solutions:
        $\sqrt a$ and $-\sqrt a$;
  \item if $a = 0$, the only solution is $0$;
  \item if $a < 0$, there is no solution.
\end{enumerate}
\end{theorem}

\begin{proof}
Rewrite $x^2 = a$ as $x^2 - \left(\sqrt a\right)^2 = 0$ when $a \geq 0$,
and factor as a \omterm{ex:g1:subtraction:difference}{difference} of squares:
\[
\left(x - \sqrt a\right)\left(x + \sqrt a\right) = 0 ,
\]
so $x = \sqrt a$ or $x = -\sqrt a$ (these coincide when $a = 0$). When
$a < 0$ there is no solution, since $x^2 \geq 0 > a$ for every $x$.
\end{proof}

\begin{omfigure}
\begin{tikzpicture}
\begin{axis}[omaxis,
  width=7.6cm, height=5.4cm,
  xmin=-3.1, xmax=3.1, ymin=-1.6, ymax=5.4,
  xtick={-1.732,1.732}, xticklabels={$-\sqrt3$,$\sqrt3$},
  ytick={3}, yticklabels={$3$}]
  \addplot[omDef, thick, domain=-2.25:2.25, samples=90] {x^2};
  \addplot[omThm, thick, domain=-2.9:2.9] {3};
  \fill (axis cs:-1.732,3) circle (1.5pt);
  \fill (axis cs:1.732,3) circle (1.5pt);
  \draw[dashed] (axis cs:-1.732,3) -- (axis cs:-1.732,0);
  \draw[dashed] (axis cs:1.732,3) -- (axis cs:1.732,0);
\end{axis}
\end{tikzpicture}

{\small Solving $x^2 = 3$ on the parabola $y = x^2$: the horizontal line
$y = 3$ cuts it at the two abscissas $\pm\sqrt3$.}
\end{omfigure}

\begin{example}
$x^2 = 16$: solutions $4$ and $-4$. \quad
$x^2 = 5$: solutions $\sqrt5$ and $-\sqrt5$ (exact values; $\sqrt5
\approx 2.236$). \quad
$x^2 = -9$: no solution. \quad
$3x^2 = 21$: divide by 3 first, $x^2 = 7$, solutions $\pm\sqrt7$.
\end{example}

\section{Exercises}

\begin{exercise}[$\star$]\label{exo:g9:sqrt:1}
Compute without a calculator:
\[
\sqrt{64}, \qquad \sqrt{100}, \qquad \sqrt{0.09}, \qquad
\sqrt{\tfrac{1}{25}}, \qquad \left(\sqrt{13}\right)^2, \qquad
\sqrt{(-4)^2} .
\]
\end{exercise}

\begin{exercise}[$\star$]\label{exo:g9:sqrt:2}
Simplify:
\[
\sqrt{50}, \qquad \sqrt{45}, \qquad \sqrt{98}, \qquad \sqrt{300} .
\]
\end{exercise}

\begin{exercise}[$\star$]\label{exo:g9:sqrt:3}
Compute and simplify:
\[
\sqrt2 \times \sqrt{18}, \qquad
\frac{\sqrt{75}}{\sqrt3}, \qquad
\sqrt5 \times \sqrt{20}, \qquad
\frac{\sqrt{8}}{\sqrt{50}} .
\]
\end{exercise}

\begin{exercise}[$\star$]\label{exo:g9:sqrt:4}
Solve: $x^2 = 36$; \quad $x^2 = 11$; \quad $x^2 + 4 = 0$; \quad
$2x^2 = 50$.
\end{exercise}

\begin{exercise}[$\star$]\label{exo:g9:sqrt:5}
Reduce to a single term: $\sqrt{20} + \sqrt{45}$, then
$3\sqrt8 - \sqrt{18} + \sqrt2$.
\end{exercise}

\begin{exercise}[$\star\star$]\label{exo:g9:sqrt:6}
Expand and simplify:
\[
\left(\sqrt3 + 1\right)^2, \qquad
\left(2\sqrt5 - 3\right)^2, \qquad
\left(\sqrt6 + \sqrt2\right)\left(\sqrt6 - \sqrt2\right).
\]
\end{exercise}

\begin{exercise}[$\star\star$]\label{exo:g9:sqrt:7}
A square field has \omterm{def:g6:measure:area}{area} $6400$ m$^2$. What is the length of its side? Of
its diagonal (exact value, then rounded to the nearest meter)?
\end{exercise}

\begin{exercise}[$\star\star$]\label{exo:g9:sqrt:8}
Show that $\dfrac{1}{\sqrt2} = \dfrac{\sqrt2}{2}$ (multiply \omterm{def:g4:fractions:def}{numerator} and
\omterm{def:g4:fractions:def}{denominator} by $\sqrt2$). Use the same trick to write $\dfrac{6}{\sqrt3}$
and $\dfrac{10}{\sqrt5}$ without a \omterm{def:g9:sqrt:def}{square root} in the \omterm{def:g4:fractions:def}{denominator}.
\end{exercise}

\begin{exercise}[$\star\star$]\label{exo:g9:sqrt:9}
True or false? Justify with a proof or a counterexample.
\begin{enumerate}
  \item For all $a, b \geq 0$: $\sqrt{ab} = \sqrt a \sqrt b$.
  \item For all $a, b \geq 0$: $\sqrt{a+b} = \sqrt a + \sqrt b$.
  \item For all $x$: $\sqrt{x^2} = x$.
  \item $\left(3\sqrt2\right)^2 = 18$.
\end{enumerate}
\end{exercise}

\begin{exercise}[$\star\star\star$]\label{exo:g9:sqrt:10}
Let $x = \sqrt{7 + 4\sqrt3}$. Compute
$\left(2 + \sqrt3\right)^2$, and deduce a simpler expression of $x$.
Same question for $\sqrt{9 - 4\sqrt5}$ (aim for a square of the form
$\left(\sqrt5 - 2\right)^2$, and mind the sign).
\end{exercise}

\section{Problem: The number that is not a fraction}

\begin{problem}[{Weekend problem --- the legendary proof that
$\sqrt2$ is irrational, its many cousins, and Heron's uncannily
good recipe for approximating it}]
\label{pb:g9:sqrt:1}
The diagonal of a unit square is a perfectly real length --- you
drew it in \cref{ex:g8:pythagoras:diagonal} --- and yet, as the
Pythagoreans discovered to their horror some twenty-five
centuries ago, \emph{no \omterm{def:g9:fractions:fraction}{fraction} whatsoever} measures it. Legend
claims the discovery was punished by drowning. This problem walks
you through the immortal proof (four questions and it is yours
for life), multiplies the victims, and ends with the recipe
engineers used for two thousand years to tame the untameable
number.

\medskip
\noindent\textbf{Part I --- The proof.}
\begin{enumerate}
  \item First hunt it down: check that
        $1.4 < \sqrt2 < 1.5$, then that $1.41 < \sqrt2 < 1.42$,
        by squaring the bounds. One more digit: between which
        three-decimal numbers does $\sqrt2$ lie?
  \item Now suppose --- for the sake of contradiction --- that
        $\sqrt2 = \frac ab$ for some \omterm{def:g9:fractions:fraction}{fraction} \emph{in lowest
        terms} (\cref{met:g9:arith:simplify}). Square both sides
        and show that $a^2 = 2 b^2$. What is the parity of
        $a^2$?
  \item The parity facts of \cref{pb:g8:pythagoras:1} say: \omterm{def:g3:numbers:evenodd}{odd
        numbers} have \omterm{def:g3:numbers:evenodd}{odd} squares. Deduce that $a$ is \omterm{def:g3:numbers:evenodd}{even}, write
        $a = 2k$, substitute --- and show that $b$ must be \omterm{def:g3:numbers:evenodd}{even}
        too.
  \item Where is the contradiction? Conclude, and state the
        theorem in full: \emph{$\sqrt2$ is irrational --- it is
        no \omterm{def:g3:division:remainder}{quotient} of whole numbers}.
  \item The proof leaned once, discreetly, on the words ``in
        lowest terms''. Point to the exact step that would
        collapse without them, and explain why assuming lowest
        terms was legitimate in the first place.
\end{enumerate}

\medskip
\noindent\textbf{Part II --- The victims multiply.}
\begin{enumerate}[resume]
  \item A second proof style, via prime factorizations
        (\cref{thm:g9:arith:factorization}): in $a^2 = 3b^2$,
        compare the parity of the \emph{\omterm{def:g8:powers:def}{exponent} of $3$} on each
        side (\cref{pb:g9:arith:1}: squares carry \omterm{def:g3:numbers:evenodd}{even}
        \omterm{def:g8:powers:def}{exponents}). Conclude that $\sqrt3$ is irrational.
  \item Run the same \omterm{def:g8:powers:def}{exponent} argument on $a^2 = n b^2$ for a
        general whole number $n$: for which $n$ does it produce
        a contradiction, and for which does it fail? State the
        complete result: $\sqrt n$ is irrational exactly
        when\,\dots
  \item Prove the shield lemma: a nonzero rational times an
        irrational is irrational. (Suppose $r x = q$ with $r, q$
        rational, $r \neq 0$, and solve for $x$.) Deduce that
        $2\sqrt2$ and $\frac{\sqrt2}{2}$ are irrational.
  \item Prove that $1 + \sqrt2$ is irrational. Then the pretty
        one: supposing $s = \sqrt2 + \sqrt3$ were rational,
        compute $s^2$, isolate $\sqrt6$, and find the
        contradiction.
  \item Temper the enthusiasm: give two irrational numbers whose
        \emph{\omterm{def:g1:addition:def}{sum}} is rational, and two whose \emph{\omterm{def:g2:mult:def}{product}} is
        rational. (Irrationality is not preserved by arithmetic
        --- each case needs its own proof.)
\end{enumerate}

\medskip
\noindent\textbf{Part III --- Heron's recipe.}
Two thousand years before calculators, Heron of Alexandria
approximated $\sqrt2$ like this: \emph{guess $x$; replace the
guess by the average of $x$ and $\frac2x$; repeat.}
\begin{enumerate}[resume]
  \item Start from the guess $x = 1$ and compute the next two
        guesses as exact \omterm{def:g9:fractions:fraction}{fractions}.
  \item Compute the third guess, again as an exact \omterm{def:g9:fractions:fraction}{fraction}, and
        its decimal value. Compare with question 1: how many
        decimals of $\sqrt2$ are already correct?
  \item The idea behind the recipe: a rectangle of \omterm{def:g6:measure:area}{area} $2$ with
        one side $x$ has the other side $\frac2x$. Show that
        $\sqrt2$ always lies \emph{between} $x$ and $\frac2x$
        (consider the two cases $x^2 > 2$ and $x^2 < 2$), so
        averaging the two sides squeezes the rectangle towards
        the square.
  \item The guesses $\frac32$, $\frac{17}{12}$,
        $\frac{577}{408}$ hide a gem: compute $3^2 - 2 \times
        2^2$, then $17^2 - 2 \times 12^2$, then
        $577^2 - 2 \times 408^2$. Part I proved
        $a^2 - 2b^2 = 0$ impossible --- how close do Heron's
        \omterm{def:g9:fractions:fraction}{fractions} come to the impossible?
  \item Finale, in two sentences: the diagonal of the unit
        square exists on paper, no \omterm{def:g9:fractions:fraction}{fraction} measures it, and its
        decimal writing can never repeat
        (\cref{pb:g9:fractions:1}). What kind of ``new numbers''
        must the number line therefore contain --- and where in
        this series is their full story told?
\end{enumerate}
\end{problem}
